\section{附录B　关键命题的证明纲要}

\subsection{B.1 Weyl 对无公共本征向量}

设 $\alpha\notin\QQ$ 且 $U_{\mathrm{scan}}V=e^{2\pi i\alpha}VU_{\mathrm{scan}}$。若存在非零 $\psi$ 同时为两者本征向量，则
\[
U_{\mathrm{scan}}V\psi=uv\psi,\qquad
VU_{\mathrm{scan}}\psi=vu\psi.
\]
由 Weyl 关系得 $uv=e^{2\pi i\alpha}uv$，推出 $e^{2\pi i\alpha}=1$，矛盾。故无公共本征向量。

\subsection{B.2 \texorpdfstring{$|X_m|=F_{m+2}$}{|X\_m|=F\_{m+2}}}

对无相邻 $1$ 的长度 $m$ 二进制词，按首位分类：若首位为 $0$，余下为任意 $X_{m-1}$；若首位为 $1$，次位必须为 $0$，余下为任意 $X_{m-2}$。故
\[
|X_m|=|X_{m-1}|+|X_{m-2}|,
\]
配合初值 $|X_1|=2$，$|X_2|=3$ 得 $|X_m|=F_{m+2}$。

\subsection{B.3 折叠平均退化度与指数标度}

由满射 $\Fold_m:\{0,1\}^m\to X_m$ 得
\[
d_m=\frac{1}{|X_m|}\sum_{w\in X_m}\big|\Fold_m^{-1}(w)\big|=\frac{2^m}{|X_m|}=\frac{2^m}{F_{m+2}}.
\]
由 Fibonacci 渐近 $F_{m+2}\asymp \varphi^{m}$ 得 $d_m\asymp (2/\varphi)^m$。

